% Irodov VI §6.3 original 159--166 (end of §6.3).
% Printed p250, user-provided PDF p259, answers PDF p355.
\atomictermsproblem{159}{250}
原子某一\emph{单重谱项}在弱磁场
$B=3.0\,\mathrm{kG}$中全劈裂能宽，
本版原印为$\Delta E=10^{4}\,\mu\mathrm{eV}$。
求其光谱项符号，并核对原书书末所列结果。
\exsolution
单重态$S=0$，所以$J=L$且$g_J=1$，
全劈裂能宽为
\[
 \Delta E=2J\mu_BB=2L\mu_BB.
\]
由于$B=0.300\,\mathrm T$，
$\mu_B B\simeq17.365\,\mu\mathrm{eV}$。
\emph{逐字按题面}$10^4\,\mu\mathrm{eV}$
代入得
\[
 L=\frac{10^4}{2\times17.365}\simeq288,
\]
即
\[
 \boxed{L\simeq288\quad
 (\text{不对应题目通常涉及的低轨道原子项})}.
\]
原书书末答案写为\({}^{1}F_3\)；
这一答案要求$L=3$，
相应全宽却是
\[
 \boxed{\Delta E_{L=3}=6\mu_BB
 \simeq104.2\,\mu\mathrm{eV}}.
\]
\excheck 原印$10^{4}\,\mu\mathrm{eV}$
和书末\({}^{1}F_3\)之间差约96倍，
可能是排版将$104$错误设成$10^4$，
但本稿不擅自将原题改为$104\,\mu\mathrm{eV}$。
此题标记为\emph{原书印刷差异待核验}。

\atomictermsproblem{160}{250}
某一由单重项间跃迁形成的光谱线
波长$\lambda=612\,\mathrm{nm}$。
在$B=10.0\,\mathrm{kG}$的磁场中，
求最外侧两塞曼分量的波长间隔。
\exsolution
单重项$g=1$，正常塞曼效应中
$\Delta m=0,\pm1$对应频移
$0,\pm\mu_B B/\hbar$。
两端的角频率总差
$\Delta\omega_{\rm ext}=2\mu_BB/\hbar$。
小频移满足
$|\Delta\lambda|=(\lambda^2/2\pi c)|\Delta\omega|$，
所以
\[
 \boxed{\Delta\lambda_{\rm ext}=
 \frac{\lambda^2\mu_B B}{\pi\hbar c}
 =\frac{eB\lambda^2}{2\pi m_ec}
 \simeq35\,\mathrm{pm}}.
\]
\excheck 题目问最外侧两条之\emph{总间隔}，
不是中央至侧线的间隔。

\atomictermsproblem{161}{250}
光谱仪分辨本领$\lambda/\delta\lambda=10^5$，
所观察单重项跃迁谱线$\lambda=536\,\mathrm{nm}$，
观察方向垂直于磁场。
求刚能分辨塞曼线分量的最小$B$。
\exsolution
光谱仪可分辨的邻线间距
$\delta\lambda=\lambda/10^5$。
正常塞曼三线的\emph{相邻}线距
\[
 \Delta\lambda_{\rm adj}
 =\frac{\lambda^2}{2\pi c}
 \frac{\mu_BB}{\hbar}
 =\frac{eB\lambda^2}{4\pi m_ec}.
\]
令$\Delta\lambda_{\rm adj}=\delta\lambda$，
\[
 \boxed{B_{\min}=
 \frac{4\pi m_ec}{e\lambda(10^5)}
 \simeq0.40\,\mathrm T=4.0\,\mathrm{kG}}.
\]
\excheck 沿垂直磁场方向可见$\pi$及$\sigma$分量，
而沿磁场方向$\pi$分量不可观察；
本题采用前者的邻线分辨条件。

\atomictermsproblem{162}{250}
在弱磁场中，
\({}^{3}D_1\to{}^{3}P_0\)谱线发生塞曼劈裂。
垂直磁场观察，相邻谱线角频率间隔
$\Delta\omega=1.32\times10^{10}\,\mathrm{s^{-1}}$。
求磁感应强度$B$。
\exsolution
上态$L=2,S=1,J=1$，
朗德因子$g_u=1/2$；
下态$J_l=0$，不劈裂。
允许跃迁的上态$m_u=0,\pm1$，
相邻频差为$g_u\mu_BB/\hbar$，
因此
\[
 \boxed{B=\frac{\hbar\Delta\omega}{g_u\mu_B}
 =\frac{2\hbar(1.32\times10^{10})}{\mu_B}
 \simeq0.300\,\mathrm T=3.0\,\mathrm{kG}}.
\]
\excheck 要用上态$g_u=1/2$，
不是盲目套普通三重线$g=1$。

\atomictermsproblem{163}{250}
钠黄线双线波长分别为
$589.59$和$589.00\,\mathrm{nm}$。
(a)求弱场下\({}^{2}P_{3/2}\)与
\({}^{2}P_{1/2}\)两项相邻塞曼子级的能距比；
(b)当\({}^{2}P_{3/2}\)的相邻磁子能级间距
为自然精细分裂的$1/\eta$，$\eta=50$时，
求磁场$B$。
\exsolution
(a) 对$L=1,S=1/2$，
朗德因子分别
$g_{3/2}=4/3$、$g_{1/2}=2/3$，
相邻塞曼间距是$g\mu_BB$，
\[
 \boxed{\frac{\delta E_{3/2}}{\delta E_{1/2}}
 =\frac{4/3}{2/3}=2:1}.
\]
(b) 自然精细分裂
\[
 \Delta E_{\rm fs}=hc\left|
 \frac1{589.00\,\mathrm{nm}}-
 \frac1{589.59\,\mathrm{nm}}\right|.
\]
令
$g_{3/2}\mu_BB=\Delta E_{\rm fs}/50$，
得到
\[
 \boxed{B=\frac{\Delta E_{\rm fs}}
 {50(4/3)\mu_B}
 \simeq0.55\,\mathrm T=5.5\,\mathrm{kG}}.
\]
\excheck (a)的比为2而非钠两种波长的比，
因为决定邻级能距的是朗德因子。

\atomictermsproblem{164}{250}
绘出弱磁场中\({}^{2}P_{3/2}\to{}^{2}S_{1/2}\)
允许跃迁关系，并求$B=4.5\,\mathrm{kG}$下
全部塞曼分量的角频率偏移。
\exsolution
上态$J_u=3/2,g_u=4/3$，
$m_u=\pm3/2,\pm1/2$；
下态$J_l=1/2,g_l=2$，
$m_l=\pm1/2$。
只保留$\Delta m=m_u-m_l=0,\pm1$的六条：
\[
 \begin{array}{c|c|c}
 m_u&m_l&g_um_u-g_lm_l\\\hline
 +3/2&+1/2&+1\\
 +1/2&+1/2&-1/3\\
 +1/2&-1/2&+5/3\\
 -1/2&+1/2&-5/3\\
 -1/2&-1/2&+1/3\\
 -3/2&-1/2&-1
 \end{array}
\]
每条偏移为
$\Delta\omega=(g_um_u-g_lm_l)\mu_BB/\hbar$。
取$B=0.45\,\mathrm T$，
$\mu_BB/\hbar\simeq3.96\times10^{10}\,\mathrm{s^{-1}}$，
因此
\[
 \boxed{\Delta\omega=
 (\pm1.32,\ \pm3.96,\ \pm6.60)\times10^{10}\,
 \mathrm{rad/s}}.
\]
\excheck 六条分量包括$\Delta m=0$两条$\pi$线，
和$\Delta m=\pm1$四条$\sigma$线；
完整图的箭头关系由表逐行确定。

\atomictermsproblem{165}{250}
图6.9有磁极N、S及小反射镜M：
从方向1（与磁场垂直）及经镜面从方向2
（等效沿磁场方向）观察同一条反常塞曼谱线。
两方向分别能看到多少条分量：
(a)\({}^{2}P_{3/2}\to{}^{2}S_{1/2}\)；
(b)\({}^{3}P_2\to{}^{3}S_1\)？
\exsolution
塞曼电偶极选择规则$\Delta m=0,\pm1$。
垂直于外场方向观察时，
$\pi$（$\Delta m=0$）与$\sigma$（$\Delta m=\pm1$）
两类分量都能出现；
沿磁场方向观察时$\pi$分量不可见。
(a) 题164已经枚举全部六条允许跃迁，
其中两条$\Delta m=0$、四条$\Delta m=\pm1$，故
\[
 \boxed{\text{(a)方向1：6条；方向2：4条}}.
\]
(b) 上态$J=2$的$m=-2,-1,0,1,2$，
下态$J=1$的$m=-1,0,1$。
枚举$\Delta m=0$有3条，
$\Delta m=+1$有3条，
$\Delta m=-1$有3条，合计9条；
沿场方向只看见六条$\sigma$分量，故
\[
 \boxed{\text{(b)方向1：9条；方向2：6条}}.
\]
\excheck 源图6.9的镜面使方向2等效观察沿磁场传播的光；
不能简单因“镜面左右对称”而认为两方向分量数相同。

\atomictermsproblem{166}{250}
弱磁场$B=3.4\,\mathrm{kG}$下，
\({}^{3}D_3\to{}^{3}P_2\)
谱线的总塞曼角频率分裂宽度$\Delta\omega$为多少？
\exsolution
上态$S=1,L=2,J=3$，
$g_u=1+[12+2-6]/24=4/3$；
下态$S=1,L=1,J=2$，
$g_l=1+[6+2-2]/12=3/2$。
频移是
\[
 \delta\omega=(g_um_u-g_lm_l)\frac{\mu_BB}{\hbar},
 \quad |m_u-m_l|\le1.
\]
枚举允许磁量子数，
最大正频移为
$(5/3)\mu_BB/\hbar$，
最小为负对称值
$-(5/3)\mu_BB/\hbar$，
故端到端角频率全宽
\[
 \boxed{\Delta\omega=
 \frac{10}{3}\frac{\mu_BB}{\hbar}
 \simeq9.97\times10^{10}\,\mathrm{rad/s}
 \approx1.0\times10^{11}\,\mathrm{rad/s}}.
\]
\excheck “总分裂”指最外两条塞曼分量的角频率差，
不能只报最大单侧位移$5\mu_BB/(3\hbar)$。
